% ==================================================================
%  TechnicalOverview.tex  --  subsection of Section 1 (Introduction)
%  The model setting (formerly Section 2.4) is placed here; its label
%  sec:model therefore points to this subsection.
%  Cross-references used here:
%    \label{sec:client} \label{sec:post}      (Section 3)
%    \label{sec:security}                     (Section 5)
%    \label{tab:notation}                     (Section 2)
% ==================================================================
\subsection{Technical Overview}
\label{sec:overview}
\label{sec:model}

\begin{figure*}[t]
  \centering
  \input{\dirI/Figures/fig_overview}
  \caption{Overview of the proposed protocol. The client sends its
    ciphertexts once and receives a single ciphertext once; every step in
    the shaded region is evaluated by the server on CKKS ciphertexts, and the
    server never obtains the client's secret key $\mathsf{sk}$.}
  \label{fig:overview}
\end{figure*}

Fig.~\ref{fig:overview} summarizes the protocol. A client holds a dataset
$X$ of $n$ points in $d$ dimensions together with the two DBSCAN
parameters, a radius $\varepsilon$ and a threshold $\mathit{MinPts}$
(Table~\ref{tab:notation}). The client generates the CKKS keys, keeps the
secret key $\mathsf{sk}$, and preprocesses its data
(Section~\ref{sec:client}). It uploads the data only once, as ciphertexts
that encrypt it, together with the public key $\mathsf{pk}$ and the
evaluation keys $\mathsf{evk}$, and may then stay offline. The client
returns only to download a single ciphertext that encrypts the cluster
labels, which it decrypts and post-processes (Section~\ref{sec:post}). A
server performs the clustering and holds no data of its own.

The server evaluates DBSCAN in the three stages shown in
Fig.~\ref{fig:overview}. \emph{Neighbor Check} decides whether the
distance between two points is at most $\varepsilon$; if so, the two
points are \emph{neighbors}. \emph{Core Check} decides whether the number
of neighbors of a point, counting the point itself, is at least
$\mathit{MinPts}$; if so, the point is a \emph{core point}. A core point
is the basic unit from which a cluster grows: the point and its neighbors
alone are already dense enough to form a cluster. \emph{Label
Propagation} then spreads labels from core points to their neighbors
until every point carries its final cluster label. All three stages run on
ciphertexts without decrypting any intermediate value, and every
comparison among them is evaluated with a composite minimax polynomial
approximation of the sign function~\cite{lee2021minimax}.

\subsubsection{Non-interactive Execution}
The protocol consists of one upload and one download. The server
evaluates the whole of DBSCAN on ciphertexts. It never sends an
intermediate value back for decryption, and it never needs the client to
be online while it computes. No other party takes part: there is no
second server with which to share the computation, and no comparison is
delegated to the client.

\subsubsection{Adversary and Public Parameters}
The server is honest-but-curious. Since the computation is carried out by
one server alone, no assumption about non-collusion is needed. Everything
the server receives is a ciphertext, a key, or one of the public
parameters $n$, $d$, $\varepsilon_n$, $\mathit{MinPts}$, $k_{\max}$, and
$R$, and the sequence of homomorphic operations it performs is determined
by these parameters alone. These few integers and the normalized radius
do not let the server infer the distribution of the data;
Section~\ref{sec:security} analyzes what each of them reveals.